% ============================================================
% Section 31: Dirac-Kirkwood 分布与密度矩阵重构
% ============================================================
\section{量子力学：Dirac--Kirkwood 分布——从联合概率到密度矩阵重构}

\begin{modelbox}{题目}
量子系统状态用密度矩阵 $\rho$ 描述（$N$ 维 Hilbert 空间）。两个非简并力学量：
\[
  A=\sum_{j=1}^N a_j|\alpha_j\rangle\langle\alpha_j|,
  \quad B=\sum_{k=1}^N b_k|\beta_k\rangle\langle\beta_k|.
\]
投影算子 $\Pi_j^A=|\alpha_j\rangle\langle\alpha_j|$，$\Pi_k^B=|\beta_k\rangle\langle\beta_k|$。

\textbf{(1)} 证明标准联合概率的两种表达式等价：
\[
  p(a_j,b_k)=p(a_j)p(b_k|a_j)
  \quad\Longleftrightarrow\quad
  p(a_j,b_k)=\operatorname{Tr}(\rho\,\Pi_j^A\Pi_k^B\Pi_j^A).
\]
\textbf{(2)} Dirac--Kirkwood 概率定义为 $q(a_j,b_k)=\operatorname{Tr}(\rho\,\Pi_j^A\Pi_k^B)$。用 $q(a_j,b_k)$ 表示密度矩阵在 $B$ 表象中的矩阵元 $\langle\beta_m|\rho|\beta_n\rangle$。
\end{modelbox}

\vspace{0.3em}
\begin{insightbox}{结论预告}
DK 概率 $q(a_j,b_k)$ 一般是复数，但包含足够信息重构密度矩阵。最终结果为
\[
  \langle\beta_m|\rho|\beta_n\rangle = \sum_{j=1}^N \frac{q(a_j,b_m)}{\langle\alpha_j|\beta_m\rangle}\,\langle\alpha_j|\beta_n\rangle.
\]
这是量子态层析（quantum tomography）的形式化基础。
\end{insightbox}


\subsection{联合概率的两种表达式等价}

\subsubsection{标准联合概率}

标准顺序测量（先 $A$ 后 $B$）的联合概率：
\begin{align*}
  p(a_j) &= \langle\alpha_j|\rho|\alpha_j\rangle = \operatorname{Tr}(\rho\,\Pi_j^A), \\
  p(b_k|a_j) &= |\langle\alpha_j|\beta_k\rangle|^2 = \langle\alpha_j|\beta_k\rangle\langle\beta_k|\alpha_j\rangle = \operatorname{Tr}(\Pi_j^A\Pi_k^B).
\end{align*}

因此：
\begin{formulabox}{标准联合概率}
\[
  p(a_j,b_k) = \langle\alpha_j|\rho|\alpha_j\rangle\,|\langle\alpha_j|\beta_k\rangle|^2.
\]
\end{formulabox}

\subsubsection{迹形式}

\begin{derivationbox}{利用迹的循环性}
考虑 $p(a_j,b_k)=\operatorname{Tr}(\rho\,\Pi_j^A\Pi_k^B\Pi_j^A)$。展开投影算符：
\begin{align*}
  \operatorname{Tr}(\rho\,\Pi_j^A\Pi_k^B\Pi_j^A)
  &= \operatorname{Tr}\bigl(\rho\,|\alpha_j\rangle\langle\alpha_j|\beta_k\rangle\langle\beta_k|\alpha_j\rangle\langle\alpha_j|\bigr) \\
  &= \langle\alpha_j|\beta_k\rangle\langle\beta_k|\alpha_j\rangle\,
     \operatorname{Tr}(\rho\,|\alpha_j\rangle\langle\alpha_j|) \\
  &= |\langle\alpha_j|\beta_k\rangle|^2\,\langle\alpha_j|\rho|\alpha_j\rangle.
\end{align*}
（第二步把 c 数 $\langle\alpha_j|\beta_k\rangle\langle\beta_k|\alpha_j\rangle$ 提出迹外；第三步用 $\operatorname{Tr}(\rho|\alpha_j\rangle\langle\alpha_j|)=\langle\alpha_j|\rho|\alpha_j\rangle$。）
\end{derivationbox}

\begin{formulabox}{等价性}
\[
  \boxed{\operatorname{Tr}(\rho\,\Pi_j^A\Pi_k^B\Pi_j^A) = \langle\alpha_j|\rho|\alpha_j\rangle\,|\langle\alpha_j|\beta_k\rangle|^2 = p(a_j)p(b_k|a_j)}
\]
\end{formulabox}

\begin{tipbox}{为什么需要两个 $\Pi_j^A$？}
中间如果只写一个 $\Pi_j^A$，即 $\operatorname{Tr}(\rho\Pi_j^A\Pi_k^B)$，结果等于 $\langle\beta_k|\rho|\alpha_j\rangle\langle\alpha_j|\beta_k\rangle$，这不是实数，更不能解释为概率。两个投影算符的``夹心''结构保证了厄米性和正定性。
\end{tipbox}


\subsection{DK 概率与密度矩阵重构}

\subsubsection{DK 概率的定义}

\begin{modelbox}{Dirac--Kirkwood 概率}
\[
  q(a_j,b_k) \equiv \operatorname{Tr}(\rho\,\Pi_j^A\Pi_k^B) = \langle\beta_k|\rho|\alpha_j\rangle\,\langle\alpha_j|\beta_k\rangle.
\]
由于 $\Pi_j^A\Pi_k^B$ 非厄米，$q(a_j,b_k)$ 一般是\textbf{复数}。
\end{modelbox}

\subsubsection{从 DK 概率提取跃迁矩阵元}

假设 $\langle\alpha_j|\beta_k\rangle\neq 0$（非简并且表象非正交时通常满足），由定义反解：
\[
  \langle\beta_k|\rho|\alpha_j\rangle = \frac{q(a_j,b_k)}{\langle\alpha_j|\beta_k\rangle}. \tag{1}
\]

\subsubsection{重构 $\langle\beta_m|\rho|\beta_n\rangle$}

\begin{derivationbox}{双重插入完备性}
在 $\langle\beta_m|\rho|\beta_n\rangle$ 两侧插入 $A$ 表象的完备性关系 $\sum_j|\alpha_j\rangle\langle\alpha_j|=I$：
\[
  \langle\beta_m|\rho|\beta_n\rangle = \sum_{i,j=1}^N \langle\beta_m|\alpha_i\rangle\langle\alpha_i|\rho|\alpha_j\rangle\langle\alpha_j|\beta_n\rangle.
\]

对 $\langle\alpha_i|\rho|\alpha_j\rangle$，再在中间插入 $B$ 表象完备性 $\sum_k|\beta_k\rangle\langle\beta_k|=I$：
\begin{align*}
  \langle\alpha_i|\rho|\alpha_j\rangle
  &= \sum_{k=1}^N \langle\alpha_i|\beta_k\rangle\langle\beta_k|\rho|\alpha_j\rangle \\
  &= \sum_{k=1}^N \langle\alpha_i|\beta_k\rangle\,
     \frac{q(a_j,b_k)}{\langle\alpha_j|\beta_k\rangle}. \quad\text{（由式 (1)）}
\end{align*}

代回原式：
\[
  \langle\beta_m|\rho|\beta_n\rangle = \sum_{i,j,k=1}^N
  \langle\beta_m|\alpha_i\rangle\langle\alpha_i|\beta_k\rangle\,
  \frac{q(a_j,b_k)}{\langle\alpha_j|\beta_k\rangle}\,
  \langle\alpha_j|\beta_n\rangle.
\]

利用 $\sum_i|\alpha_i\rangle\langle\alpha_i|=I$ 化简 $i$ 求和：
\[
  \sum_{i=1}^N \langle\beta_m|\alpha_i\rangle\langle\alpha_i|\beta_k\rangle
  = \langle\beta_m|\beta_k\rangle = \delta_{mk}.
\]

因此 $k$ 求和坍缩为 $k=m$：
\begin{align*}
  \langle\beta_m|\rho|\beta_n\rangle
  &= \sum_{j=1}^N \frac{q(a_j,b_m)}{\langle\alpha_j|\beta_m\rangle}\,
     \langle\alpha_j|\beta_n\rangle.
\end{align*}
\end{derivationbox}

\begin{formulabox}{密度矩阵元的 DK 表示}
\[
  \boxed{\langle\beta_m|\rho|\beta_n\rangle = \sum_{j=1}^N \frac{q(a_j,b_m)}{\langle\alpha_j|\beta_m\rangle}\,\langle\alpha_j|\beta_n\rangle}
\]
其中 $q(a_j,b_m)=\operatorname{Tr}(\rho\,\Pi_j^A\Pi_m^B)$，要求分母 $\langle\alpha_j|\beta_m\rangle\neq 0$。
\end{formulabox}

\begin{insightbox}{$m=n$ 时的对角元}
当 $m=n$，上式给出 $B$ 表象中的概率分布：
\[
  \langle\beta_m|\rho|\beta_m\rangle = \sum_{j=1}^N q(a_j,b_m).
\]
这表明虽然单个 $q(a_j,b_k)$ 可复，但对 $j$ 求和后恢复为实的、非负的概率。这与复函数实部/虚部的约束结构类似。
\end{insightbox}


\subsection{物理意义与拓展}

\subsubsection{DK 分布 vs 标准概率}

\begin{insightbox}{两种概率的本质区别}
\begin{center}
\begin{tabular}{lcc}
\toprule
\textbf{性质} & \textbf{标准联合概率 $p(a_j,b_k)$} & \textbf{DK 概率 $q(a_j,b_k)$} \\
\midrule
数值 & 实数，$\ge 0$ & 一般为复数 \\
厄米性 & $\Pi_j^A\Pi_k^B\Pi_j^A$ 厄米 & $\Pi_j^A\Pi_k^B$ 非厄米 \\
概率公理 & 满足 Kolmogorov 公理 & 不满足 \\
信息内容 & 仅经典关联 & 包含量子相干 \\
密度矩阵重构 & 不充分（丢失相位） & 充分（含相位信息） \\
\bottomrule
\end{tabular}
\end{center}
\end{insightbox}

\subsubsection{与 Wigner 函数的联系}

\begin{insightbox}{准概率分布的统一框架}
DK 分布 $q(a_j,b_k)$ 是离散变量版本的\textbf{准概率分布}。在连续极限下（如 $X$ 和 $P$ 表象），DK 分布与 \textbf{Wigner 函数} $W(x,p)$ 有密切联系：
\begin{itemize}[nosep]
  \item Wigner 函数也是准概率（可负）
  \item DK 分布的复性对应 Wigner 函数的``非经典性''
  \item 两者都包含重构密度矩阵的完整信息
\end{itemize}
\textbf{历史背景}：Kirkwood (1933) 首次引入此分布；Dirac 更早讨论了类似的复概率概念。近年来在弱测量（weak measurement）和量子层析中被重新发掘。
\end{insightbox}


\subsection{常见错误与考场策略}

\begin{errorbox}{易错点}
\begin{enumerate}
  \item \textbf{迹的循环性误用}：$\operatorname{Tr}(ABC)=\operatorname{Tr}(CAB)$ 要求算符在迹下可循环，但不要把 c 数与算符混淆。如 $|\alpha_j\rangle\langle\alpha_j|\beta_k\rangle\langle\beta_k|$ 中，$\langle\alpha_j|\beta_k\rangle$ 是 c 数，可直接提出。
  \item \textbf{分母为零}：重构公式要求 $\langle\alpha_j|\beta_m\rangle\neq 0$。若 $A$ 和 $B$ 有共同本征态（即 $[A,B]=0$），则某些 $\langle\alpha_j|\beta_m\rangle=0$ 或 $1$，此时需要特殊处理。
  \item \textbf{DK 概率的复数性}：$q(a_j,b_k)$ 不是测量概率，不能直接代入 $P=|c|^2$ 公式。其物理意义是准概率的``复振幅''。
  \item \textbf{完备性插入位置}：重构 $\langle\beta_m|\rho|\beta_n\rangle$ 时，需要在 $\rho$ 的左右两侧\textbf{分别}插入完备性，而不是只插一侧。
\end{enumerate}
\end{errorbox}

\begin{cheatbox}{DK 分布：考场速查}
\textbf{标准联合概率}：$p(a_j,b_k)=\operatorname{Tr}(\rho\,\Pi_j^A\Pi_k^B\Pi_j^A)=\langle\alpha_j|\rho|\alpha_j\rangle\,|\langle\alpha_j|\beta_k\rangle|^2$

\textbf{DK 概率}：$q(a_j,b_k)=\operatorname{Tr}(\rho\,\Pi_j^A\Pi_k^B)=\langle\beta_k|\rho|\alpha_j\rangle\langle\alpha_j|\beta_k\rangle$

\textbf{重构公式}：
\[
  \langle\beta_m|\rho|\beta_n\rangle = \sum_j \frac{q(a_j,b_m)}{\langle\alpha_j|\beta_m\rangle}\,\langle\alpha_j|\beta_n\rangle
\]
\textbf{对角元}：$\langle\beta_m|\rho|\beta_m\rangle=\sum_j q(a_j,b_m)$（实数、非负）
\end{cheatbox}
